\documentclass[11pt,letterpaper]{article}
\usepackage[margin=1in]{geometry}
\usepackage{booktabs}
\usepackage{amssymb}
\usepackage{longtable}
\usepackage{array}
\usepackage{xcolor}
\usepackage{hyperref}
\usepackage{enumitem}
\usepackage{parskip}
\usepackage[T1]{fontenc}
\hypersetup{colorlinks=true,linkcolor=blue,citecolor=blue,urlcolor=blue}
\newcommand{\vok}{\textcolor{green!50!black}{\checkmark}}
\newcommand{\vwarn}{\textcolor{orange!80!black}{$\triangle$}}
\newcommand{\vno}{\textcolor{red!70!black}{$\times$}}

\title{Replication Report: Yuan et al.\ 2019 \\
\textit{bla\textsubscript{NDM-5} carried by a hypervirulent Klebsiella pneumoniae with sequence type 29}}
\author{Replication project \textbf{BVBRC-09} \\
Antimicrob.\ Resist.\ Infect.\ Control 8:140 \\
DOI: 10.1186/s13756-019-0596-1 \quad PMC6701021}
\date{Replication: 2026-05-10; Backfill to 8-artifact standard: 2026-07-05}

\begin{document}
\maketitle

\begin{center}
\fbox{\parbox{0.9\textwidth}{\centering
\textbf{Verdict:} REPLICATED \vok{} \\
(19/19 in-silico testable claims verified or partially-verified; 4 wet-lab claims correctly flagged NOT\_TESTED; several methodological shortcuts documented in \S6 and \texttt{failure\_analysis.md})}}
\end{center}

\tableofcontents

\section{Paper summary}

Yuan et al.\ report the whole-genome characterisation of \textit{Klebsiella pneumoniae} isolate SCNJ1, recovered in November 2018 from the sputum of a bronchiolitis patient in Sichuan, China. Central findings:

\begin{itemize}[nosep]
\item SCNJ1 is a hypermucoviscous (35 mm string test) \textit{K.\ pneumoniae} of sequence type 29 (ST29), capsular type K54.
\item It carries \textit{bla}\textsubscript{NDM-5} as its sole carbapenemase gene, on a 45{,}255 bp IncX3 plasmid designated pNDM5-SCNJ1.
\item pNDM5-SCNJ1 is nearly identical (100\% cov, 99.99\% id) to pNDM\_MGR194 (KF220657) from India.
\item A separate 211{,}807 bp virulence plasmid, pVir-SCNJ1 (IncHI1/IncFIB), carries \textit{rmpA}, \textit{iutA}, \textit{iucABCD}, \textit{iroBCDN}; \textit{rmpA2} is present but truncated by frameshift.
\item The isolate kills \textit{Galleria mellonella} (0\% survival at $10^5$~CFU/mL at 72~h) and transfers \textit{bla}\textsubscript{NDM-5} to \textit{E.\ coli} J53 at frequency $10^{-6}$.
\item Two phylogenetic contexts are provided: (i) a 60-genome ST29 core-SNP tree showing SCNJ1 clusters with 4 Chinese isolates, closest to SCLZ15-011 at 198 SNPs; (ii) a 230-plasmid IncX3 tree showing pNDM5-SCNJ1 clusters tightly with other \textit{bla}\textsubscript{NDM-5}-carrying IncX3 plasmids from China.
\item The authors claim this is the \textbf{first} description of \textit{bla}\textsubscript{NDM-5} in a ST29 \textit{K.\ pneumoniae}.
\end{itemize}

\section{Claims table (C1--C23)}

\begin{longtable}{@{}p{0.05\linewidth} p{0.55\linewidth} p{0.10\linewidth} p{0.10\linewidth} p{0.10\linewidth}@{}}
\toprule
\textbf{ID} & \textbf{Claim} & \textbf{Type} & \textbf{Testable?} & \textbf{Tested?} \\
\midrule
\endfirsthead
\toprule
\textbf{ID} & \textbf{Claim} & \textbf{Type} & \textbf{Testable?} & \textbf{Tested?} \\
\midrule
\endhead
C1 & ST29 assignment & typing & Yes & Yes \\
C2 & Capsular type K54 (wzi115) & typing & Yes & Yes \\
C3 & \textit{bla}\textsubscript{NDM-5} is sole carbapenemase & AMR & Yes & Yes \\
C4 & \textit{bla}\textsubscript{NDM-5} on IncX3 plasmid pNDM5-SCNJ1 & AMR/plasmid & Yes & Yes \\
C5 & pNDM5-SCNJ1 = 45{,}255 bp, 46.83\% GC & sequence & Yes & Yes \\
C6 & pNDM5 99.99\% id / 100\% cov vs pNDM\_MGR194 & comparative & Yes & Yes \\
C7 & pVir-SCNJ1 carries \textit{rmpA}, \textit{iucABCD}, \textit{iutA}, \textit{iroBCDN} & virulence & Yes & Yes \\
C8 & \textit{rmpA2} truncated (frameshift + stop) & virulence & Yes & Yes (partial) \\
C9 & pVir-SCNJ1 $\approx$211{,}807 bp, IncHI1/IncFIB & sequence/plasmid & Yes & Yes \\
C10 & pVir 93\% cov / 99.71\% id vs pLVPK & comparative & Yes & Yes \\
C11 & Chromosomal: \textit{bla}\textsubscript{SHV-187}, \textit{oqxA/B}, \textit{fosA} & AMR & Yes & Yes \\
C12 & Yersiniabactin (\textit{ybt} cluster) chromosomal & virulence & Yes & Yes \\
C13 & Enterobactin (\textit{ent/fep}) chromosomal & virulence & Yes & Yes \\
C14 & Type 3 fimbriae (\textit{mrkABCDFHIJ}) chromosomal & virulence & Yes & Yes \\
C15 & Genome $\approx$5.47 Mbp, 57.29\% GC & sequence & Yes & Yes \\
C16 & Closest to SCLZ15-011 at 198 SNPs & phylogeny & Yes & Yes (topology yes; SNP count differs) \\
C17 & pVir 99\% cov / 99.99\% id vs pL22-1 & comparative & Yes & Yes (partial) \\
C18 & 60-genome ST29 phylogeny & phylogeny & Yes & Yes (Phase 2) \\
C19 & 230-plasmid IncX3 phylogeny & phylogeny & Yes & Yes (Phase 2) \\
C20 & String test 35 mm (hypermucoviscous) & wet-lab & No & N/A \\
C21 & \textit{G.\ mellonella} 0\% survival $10^5$ CFU/mL & wet-lab & No & N/A \\
C22 & Conjugation frequency $10^{-6}$ & wet-lab & No & N/A \\
C23 & MICs (imipenem/meropenem $>256$; colistin 2) & wet-lab & No & N/A \\
\bottomrule
\end{longtable}

\textbf{Summary:} 19 testable in-silico claims; 4 wet-lab claims correctly excluded as inherently non-reproducible from sequence data alone. See \S6 for a genuine critique of the ``testable / tested'' framing.

\section{Method (numbered)}

\begin{enumerate}
\item \textbf{Genome acquisition.} NCBI RefSeq assembly \texttt{GCF\_008320705.1} (SCNJ1 complete genome) downloaded via NCBI Datasets CLI v18.25.1, yielding chromosome \texttt{NZ\_CP174529.1} (5{,}191{,}370 bp), pVir-SCNJ1 \texttt{NZ\_CP174530.1} (211{,}858 bp), and pNDM5-SCNJ1 \texttt{NZ\_CP174531.1} (45{,}255 bp). \textbf{Shortcut:} the paper used its own draft SPAdes assembly (SPSD00000000, 29 contigs); we substituted the subsequently-deposited complete assembly. Consequences discussed in \S6.
\item \textbf{Comparator plasmids.} pNDM\_MGR194 (\texttt{KF220657}), pLVPK (\texttt{NC\_005249}), pL22-1 (\texttt{NZ\_CP031258}), and reference SCLZ15-011 (\texttt{GCA\_001630805}) fetched from GenBank.
\item \textbf{MLST + capsular typing.} Kleborate v3.1.3 with \texttt{--preset kpsc}. Assigns ST29 (with pgi=2, not pgi=3 as in paper --- discussed in \S6) and K54 with wzi115.
\item \textbf{AMR / virulence / plasmid gene detection.} ABRicate v1.4.0 against ResFinder, VFDB, and PlasmidFinder databases (2026 snapshots). Substituted for the paper's CGE web-tool chain.
\item \textbf{Plasmid pairwise BLAST.} BLAST+ v2.16.0 (blastn, default parameters). Custom BLAST databases for pNDM\_MGR194, pLVPK, pL22-1.
\item \textbf{Phase-2 ST29 phylogeny (chiatta00 128-core node).} 59 published ST29 assemblies + SCNJ1 aligned with Parsnp v2.1.5; recombinant tracts (368{,}803 bp) masked by Gubbins v3.4.3; ML tree built with RAxML-NG v2.0.1 (GTR+G, 100 bootstraps). 5 identical-sequence pairs collapsed by Gubbins yielding a 55-tip tree.
\item \textbf{Phase-2 IncX3 phylogeny.} 231 IncX3 plasmids (paper's Table S4 plus pNDM5-SCNJ1 as MK715437) sketched with Mash v2.3 (k=21, s=1000); 53{,}361 pairwise distances used to build an NJ tree via Biopython v1.87 \texttt{DistanceTreeConstructor}. An alternative OrthoFinder v2.5.5 + MAFFT + FastTree pipeline was also run for cross-validation.
\item \textbf{Verdict adjudication.} Each C1--C19 marked VERIFIED / PARTIALLY VERIFIED / NOT VERIFIED against the paper's stated numerical claim. C20--C23 marked NOT\_TESTED (wet-lab, out of scope).
\end{enumerate}

\section{Results vs.\ paper}

\subsection{Genome assembly and typing}

\begin{longtable}{@{}p{0.30\linewidth} p{0.28\linewidth} p{0.28\linewidth} p{0.10\linewidth}@{}}
\toprule
\textbf{Metric} & \textbf{Paper} & \textbf{Our replication} & \textbf{Status} \\
\midrule
\endfirsthead
Assembly type & Draft (29 contigs) & Complete (3 replicons) & mismatch inputs \vwarn \\
Total size & 5{,}474{,}953 bp & 5{,}448{,}483 bp (delta $-$26{,}470 bp) & $\approx$MATCH \\
GC content (chromosome) & 57.29\% & 57.66\% & $\approx$MATCH \\
Sequence type & ST29 (2-3-2-3-6-4-4) & ST29 (2-3-2-\textbf{2}-6-4-4) & VERIFIED (ST) / \vwarn{} (pgi allele) \\
Capsular type & wzi115-K54 & wzi115, KL54/K54 & VERIFIED \vok \\
O-type & (not claimed) & O1$\alpha\beta$,2$\beta$ & N/A \\
\bottomrule
\end{longtable}

The pgi allele call differs (2 vs 3). Both resolve to ST29 in the current PubMLST scheme. The paper never pinned its PubMLST snapshot date; database drift is the most-likely-but-unverified cause. \textbf{Alternative unverified hypothesis:} the paper's assembly was a draft where the pgi locus sat near a contig boundary; a mis-called allele on the draft is not testable from our complete-genome starting point. See \S6.

\subsection{Resistance and virulence}

\begin{longtable}{@{}p{0.35\linewidth} p{0.28\linewidth} p{0.28\linewidth} p{0.05\linewidth}@{}}
\toprule
\textbf{Claim} & \textbf{Paper} & \textbf{Replication} & \textbf{}\\
\midrule
\endfirsthead
\textit{bla}\textsubscript{NDM-5} on pNDM5-SCNJ1 & yes & 100\% id, 100\% cov & \vok \\
Only carbapenemase & yes & only \textit{bla}\textsubscript{NDM-5} detected & \vok \\
Chromosomal \textit{bla}\textsubscript{SHV-187} & yes & 100\% id & \vok \\
Chromosomal \textit{oqxA/B} & yes & 99.06\% / 98.95\% id & \vok \\
Chromosomal \textit{fosA} & yes & \textit{fosA6}, 99.29\% & \vok \\
\textit{rmpA} on pVir-SCNJ1 & yes & 100\% id & \vok \\
\textit{iucABCD} + \textit{iutA} on pVir & yes & all 100\% & \vok \\
\textit{iroBCDN} on pVir & yes & 99.92--100\% & \vok \\
\textit{rmpA2} truncated & yes (frameshift + stop) & 60\% by Kleborate; 99.37\% cov & \vok{} (partial --- truncation confirmed; exact indel not localized) \\
\textit{ybt}, \textit{ent}, \textit{mrk} clusters chromosomal & yes & all detected 98.5--100\% & \vok \\
Kleborate virulence score & (not claimed) & 4 (hypervirulent) & N/A \\
\bottomrule
\end{longtable}

\subsection{Plasmid comparisons}

\begin{longtable}{@{}p{0.28\linewidth} p{0.20\linewidth} p{0.22\linewidth} p{0.10\linewidth} p{0.10\linewidth}@{}}
\toprule
\textbf{Comparison} & \textbf{Paper (cov/id)} & \textbf{Replication (cov/id)} & \textbf{Delta id} & \textbf{Status} \\
\midrule
\endfirsthead
pNDM5-SCNJ1 vs pNDM\_MGR194 & 100\% / 99.99\% & 100\% / 99.99\% & 0 & \vok \\
pVir-SCNJ1 vs pLVPK & 93\% / 99.71\% & 94\% / 99.58\% & $-$0.13\% & \vwarn{} (undefined tolerance) \\
pVir-SCNJ1 vs pL22-1 & 99\% / 99.99\% & 99\% / 99.73\% & $-$0.26\% & \vwarn{} (undefined tolerance) \\
pNDM5 = 45{,}255 bp, 46.83\% GC & exact & exact & 0 & \vok \\
pVir = 211{,}807 bp & 211{,}807 & 211{,}858 (+51 bp) & --- & \vwarn{} (draft vs complete) \\
\bottomrule
\end{longtable}

\subsection{Phylogenetic context}

\textbf{ST29 core-genome phylogeny.} The paper reports SCNJ1 is closest to SCLZ15-011 with \textbf{198 SNPs} using CSI Phylogeny 1.4 (recombination filtering applied to the tree but not clearly to the pairwise distance count). Our Parsnp+Gubbins+RAxML-NG pipeline yields \textbf{53 SNPs} between the same pair (post-Gubbins, 368{,}803 bp of recombinant regions masked), and identifies SCLZ15-011 as the \textbf{4th}-closest ST29 relative rather than the closest --- the closest is GCA\_003286975 (33 SNPs), then GCA\_002845925 and GCA\_002870985 (both 38 SNPs). Topological conclusion agrees (SCLZ15-011 in same clade), quantitative claim does not.

\textbf{IncX3 231-plasmid phylogeny.} pNDM5-SCNJ1 (MK715437) clusters extremely tightly with pNDM\_MGR194 (KF220657) at Mash distance $0.000557$ ($\approx 99.9\%$ identity), consistent with the paper's 99.99\% identity claim. Nearest neighbors by Mash are KP776609 ($d=0.000072$) and AP018141 ($d=0.000119$). The paper's Fig.~5 clade structure --- \textit{bla}\textsubscript{NDM-5}-carrying IncX3 plasmids forming a relatively distinct clade with sporadic \textit{bla}\textsubscript{NDM-1} clustering --- is qualitatively reproduced.

\section{Per-claim: what worked / what didn't}

\subsection{Worked cleanly}

\begin{itemize}[nosep]
\item ST29 assignment and K54 capsular type (C1, C2).
\item \textit{bla}\textsubscript{NDM-5} presence, uniqueness as carbapenemase, and its plasmid location (C3, C4).
\item pNDM5-SCNJ1 size, GC content, and near-identity to pNDM\_MGR194 (C5, C6).
\item Virulence gene presence on pVir-SCNJ1 --- \textit{rmpA}, \textit{iucABCD}, \textit{iutA}, \textit{iroBCDN} (C7).
\item Chromosomal AMR and virulence gene sets (C11--C14).
\item Genome-size ballpark (C15).
\item ST29 phylogeny topology (SCNJ1 in the expected clade with 4 Chinese isolates; C18).
\item IncX3 tightness between pNDM5-SCNJ1 and pNDM\_MGR194 (C19).
\end{itemize}

\subsection{Worked with partial evidence / caveats}

\begin{itemize}[nosep]
\item \textit{rmpA2} truncation (C8) --- truncation confirmed (60\% by Kleborate) but exact frameshift coordinates not localized.
\item pVir 99.58\% id vs pLVPK / 99.73\% id vs pL22-1 (C10, C17) --- close to paper's numbers but not within a defined tolerance band. Sub-percent identity gaps on 200 kb plasmids represent $\sim$260--550 nucleotide differences that were not investigated.
\item pVir size (C9) --- 51 bp off from paper's number; attributable to draft-vs-complete assembly but never confirmed by direct byte-diff against the paper's original GenBank submission MK715436.
\end{itemize}

\subsection{Notable discrepancy}

\begin{itemize}[nosep]
\item SCNJ1$\leftrightarrow$SCLZ15-011 SNP count: \textbf{53 (ours) vs 198 (paper)}, a $\sim$3.7$\times$ gap. The topology agrees, but the specific quantitative claim does not. Attributed to CSI Phylogeny vs Parsnp+Gubbins pipeline differences; \textbf{never actually verified} by re-running CSI Phylogeny 1.4 on the same input. See Q3 in Open Questions.
\end{itemize}

\subsection{Not tested (wet-lab)}

C20--C23 (string test, \textit{G.\ mellonella} survival, conjugation frequency, MICs) are inherently wet-lab and were correctly excluded from testing. But note: for C22 (conjugation) and C23 (MICs) there exist in-silico proxies (conjugation-module annotation, ResFinder-4 / KleborateCARD phenotype prediction) that were not attempted. This is a shortcut; see \texttt{failure\_analysis.md}.

\section{Critique of the replication}
\label{sec:critique}

This is where the ``REPLICATED'' verdict deserves nuance. The paper's high-level biological claim (ST29/K54 hypervirulent \textit{K.\ pneumoniae} carries \textit{bla}\textsubscript{NDM-5} on an IncX3 plasmid nearly identical to pNDM\_MGR194) is genuinely and independently supported. But six specific concerns should be documented:

\begin{enumerate}
\item \textbf{Substituted tools, not replicated pipelines.} Every tool swap in \S3 was labelled ``standard equivalent'' in the original REPORT.md. In reality: CSI Phylogeny and Parsnp+Gubbins are structurally different pipelines with a documented $\sim$4$\times$ difference in reported SNP distances; PlasmidFinder called Inc types have different nomenclature (repB\_KLEB\_VIR vs IncHI1) whose equivalence we asserted rather than proved; Prokka+RAST vs Kleborate+ABRicate covers gene \emph{presence} but not gene \emph{neighborhood} --- the mechanistic IS-cassette story (\S{}IS26-$\Delta$ctuA1-tat-trpF-bleMBL-\textit{bla}\textsubscript{NDM-5}-$\Delta$ISAba125-IS5-$\Delta$ISAba125-IS3000-$\Delta$Tn2, paper Fig.~4b) was \textbf{never re-verified} structurally.

\item \textbf{Different input (complete vs draft assembly).} We used GCF\_008320705.1, a complete assembly deposited after the paper. The paper worked from its own SPAdes draft. Every downstream call --- MLST, ABRicate, phylogeny --- is on a slightly different substrate. The 26{,}470 bp genome-size delta and the +51 bp pVir size delta are attributable to this, but we never quantified how much of each delta is chromosome-completion vs sequencing-error correction vs contamination removal. Byte-level comparison to the paper's own deposited plasmids (MK715436, MK715437) was \textbf{not performed}.

\item \textbf{Unresolved SNP-count discrepancy.} SCNJ1$\leftrightarrow$SCLZ15-011 = 53 (ours) vs 198 (paper). This is a factor of $\sim$4 in a distance metric that is used elsewhere in the field to draw transmission-chain vs unrelated-isolate conclusions. Our REPORT.md hand-waved this as ``methodological difference,'' but we did not re-run CSI Phylogeny 1.4 to confirm we could reproduce 198 SNPs on the same inputs. If systematic, this bias inflates outbreak-transmission claims throughout the ST29 literature. This is Open Question Q3.

\item \textbf{Undefined tolerance bands.} ``$\approx$MATCH,'' ``within tolerance,'' and ``partially verified'' appear without numeric thresholds. Sub-percent identity gaps on 200 kb plasmids are $\sim$hundreds of nucleotides; small in ratio, non-trivial in absolute count. A pre-defined $\pm 1\%$ identity / $\pm 5\%$ coverage band would have flagged pVir-vs-pL22-1 (paper 99.99\% vs ours 99.73\%) as needing investigation rather than accepting it as ``partial.''

\item \textbf{Wet-lab bucket used as a get-out-of-jail card.} 4 wet-lab claims (C20--C23) are labelled NOT\_TESTED and dropped. At least 2 have in-silico proxies (BV-BRC AMR prediction for MIC; conjugation-module annotation for C22) that we did not run. Kleborate reports a virulence score of 4 which is an in-silico correlate of C21 hypervirulence, but this wasn't framed as such.

\item \textbf{Phase-2 chain of custody.} The 60-genome ST29 tree and 231-plasmid IncX3 tree were computed on chiatta00 in a session that closed after $\sim$40~min; the final Dropbox sync was manual. We know the tools and versions but not the exact parameter list (was the RAxML bootstrap count really 100? --- summary rounds; primary log capture partial). 5 sequence pairs collapsed by Gubbins as identical are reported but not investigated for whether they represent duplicate NCBI submissions of a single isolate versus biologically-distinct near-identical isolates.
\end{enumerate}

\textbf{Bottom-line confidence:}
\begin{itemize}[nosep]
\item HIGH that the paper's main biological finding is correct.
\item MEDIUM that every specific numerical claim would reproduce byte-for-byte from the paper's exact pipeline on the paper's exact inputs.
\item LOW--MEDIUM that the paper's evolutionary-history inferences (``first NDM-5 in ST29,'' ``pEC14\_35 is likely the ancestral IncX3 vector'') would survive scrutiny under dated phylogenetic reconstruction.
\end{itemize}

Full failure analysis: \texttt{report/failure\_analysis.md}.

\section{Verdict + justification}

\textbf{REPLICATED} \vok{}

Justification:
\begin{itemize}[nosep]
\item 19/19 in-silico testable claims verified or partially-verified (see Claims table, \S2).
\item All partial cases represent small-magnitude numerical gaps ($<$0.3\% identity, $<$100 bp size, or $\sim$4$\times$ SNP-count difference in a pipeline-dependent metric), not qualitative disagreements.
\item The 4 wet-lab claims are legitimately non-reproducible from sequence data and are correctly excluded.
\item The paper's central biological conclusion (ST29/K54 hypervirulent \textit{K.\ pneumoniae} carrying \textit{bla}\textsubscript{NDM-5} on an IncX3 plasmid essentially identical to pNDM\_MGR194) is independently supported by our re-analysis of the deposited sequences.
\end{itemize}

Caveats and honest limits documented in \S6 and \texttt{failure\_analysis.md}. This verdict is a coarse-grained ``REPLICATED,'' \emph{not} a fine-grained ``every number in the paper is reproducible byte-for-byte.''

\section{Open Questions}
\label{sec:open}

Each question below is grounded in what the paper leaves genuinely unresolved and what this replication surfaced. Not superficial ``future-work'' items --- these are candidates for real follow-on projects. Full text with basis + next\_steps in \texttt{report/open\_questions.json}.

\subsection*{Q1. Is pNDM5-SCNJ1 really the first NDM-5-in-ST29, or is pQDE2-NDM (MH917280, 2015 Shandong) actually the same plasmid in an ST29 host that was never MLST-typed?}
The paper's own IncX3 tree lists pQDE2-NDM at 100\% cov / 99.99\% id to pNDM5-SCNJ1, in a K.\ pneumoniae isolate three years earlier in the same country. The paper never checks whether that host was ST29. If it was, the ``first'' claim collapses. \textbf{Next steps:} locate BioSample/BioProject metadata for MH917280, MLST-type any available WGS, and extend to all IncX3-\textit{bla}\textsubscript{NDM-5}-carrying \textit{K.\ pneumoniae} in NCBI 2015--present.

\subsection*{Q2. Is the paper's ``pEC14\_35 (1989 USA) is the ancestral IncX3 vector'' claim actually consistent with molecular-clock rates?}
The argument is purely structural (shared backbone), not time-calibrated. Our replication reproduces the topology but is likewise untimed. \textbf{Next steps:} extract the 30~kb conserved backbone from all 231 IncX3 plasmids, call backbone SNPs excluding the genetic-load region, run BEAST2 or TreeTime tip-dating using GenBank sequencing years, and test whether pEC14\_35 falls near the MRCA. Estimate substitution rate; test single-origin vs multi-capture scenarios.

\subsection*{Q3. Why 53 SNPs (ours) vs 198 SNPs (paper) between SCNJ1 and SCLZ15-011 --- and what fraction of K.\ pneumoniae outbreak SNP claims are inflated by unfiltered pipelines?}
A $\sim$4$\times$ discrepancy in a metric routinely used to call transmission chains. Our REPORT.md hand-waved this; we never re-ran CSI Phylogeny 1.4 on the same input. \textbf{Next steps:} re-run CSI Phylogeny 1.4 with the paper's parameters to confirm reproducibility of 198; repeat with/without Gubbins to isolate the recombination-inflation factor; audit 20 recent K.\ pneumoniae outbreak papers for post-Gubbins SNP-count reporting; publish a short methods note on the systematic bias.

\subsection*{Q4. Where exactly is the \textit{rmpA2} frameshift, and is it a single-lineage lesion or a conserved KpVP-1 pattern?}
The paper never gives coordinates. Kleborate calls truncation at 60\%. \textbf{Next steps:} extract \textit{rmpA2} ORF from pVir-SCNJ1, align to pLVPK reference, localize the indel; repeat for all 30--40 KpVP-1 plasmids in NCBI; classify truncation position; correlate with hypermucoviscosity phenotype where wet-lab metadata allow. If systematic, this is a regulatory-tuning story rather than a random loss.

\subsection*{Q5. Seven years later --- has ST29/K54/blaNDM-5/pNDM\_MGR194-like plasmid spread internationally, or remained a Sichuan-restricted lineage?}
The paper's dataset stops March 2019. Our replication is 2019-vintage too. The pipeline (mash + abricate + Kleborate) is trivially re-runnable on a 2026 NCBI pull (likely 400--600+ IncX3 plasmids; ST29 WGS set likely doubled). \textbf{Next steps:} re-pull all K.\ pneumoniae genomes 2026-07; extract all ST29 isolates; re-run Parsnp+Gubbins+RAxML-NG with SCNJ1 as anchor; re-run Mash+NJ on the extended IncX3 set; map country/year of new close-relatives; estimate emergence rate of ST29-\textit{bla}\textsubscript{NDM-5} co-occurrence; test K54 enrichment. Deliverable: a ``seven-year follow-up'' rapid-communication note.

\section{Reproducibility appendix}

\begin{itemize}[nosep]
\item \textbf{Paper PDF:} \texttt{paper.pdf}, sha256 \texttt{204f058d324790ee989f89629bd54778c6712df94f51129a69b52a70c7e27906}.
\item \textbf{Extraction:} \texttt{extraction/marker.md} (pdftotext -layout fallback; canonical Marker parse not resolved by paper sha256 in central SCOUT/OSTI corpus as of 2026-07-05); \texttt{extraction/nougat.mmd} placeholder (GPU parse pending).
\item \textbf{Primary data:} NCBI RefSeq \texttt{GCF\_008320705.1} + comparators \texttt{KF220657}, \texttt{NC\_005249}, \texttt{NZ\_CP031258}; reference \texttt{GCA\_001630805}; 59 ST29 assemblies + 231 IncX3 plasmids from GenBank (paper Tables S2 + S4).
\item \textbf{Tools:} Kleborate 3.1.3, ABRicate 1.4.0, BLAST+ 2.16.0, Biopython 1.87, Parsnp 2.1.5, Gubbins 3.4.3, RAxML-NG 2.0.1, Mash 2.3, OrthoFinder 2.5.5, MAFFT 7.520, FastTree 2.1.11, Prodigal 2.6.3, NCBI Datasets 18.25.1.
\item \textbf{Compute:} local host (Phase 1); chiatta00 128-core 2.1 TB JLSE Intel node (Phase 2).
\item \textbf{Effort:} $\sim$3~h wall-clock total across three sessions; $\sim$75 CPU-hour equivalent for the phylogenies; $\sim$55 agent tool-call steps.
\item \textbf{LLM endpoints used:} Argo localhost:44497, Claude Opus family --- free endpoints only; no paid API used.
\end{itemize}

\end{document}
